% ================================================================
%  chapters/parts/ch11-d.tex
%  第 11 章  电化学催化研究方法（分片 d）
%  节次：§ 11.5.1  阻抗谱基本原理
%  源：S2（下册）第 11 章 §11.5.1，书内 p.855—861 / PDF p.343—349
%      （实测：该节源文本在 sources/parts/ch11-b.txt 第 864—1094 行，
%        起于「11.5.1 阻抗谱基本原理」标题，止于「11.5.2 阻抗技术在电催化
%        领域中的应用」标题之前，含原书图 11-32 至图 11-41、表 11-1，
%        公式编号（11-115）至（11-120）；源文本未给（11-121）以后的编号。）
%
%  【顺序问题——已于 2026-09-29 由装配层解决】
%  装配器 chapters/ch11.tex 原先按 a→b→c→d 的 \IfFileExists 顺序引入，
%  本分片实际排在 ch11-c（§11.5 后半与 §11.6）之后，而不是 §11.5.1 应处的
%  §11.4 之后，由此产生三处节序错乱：
%    (1) 正文顺序变成 §11.4 → §11.5.1 → §11.5.2 → §11.6 → §11.5（ch11-c）：
%        「阻抗谱基本原理」与其应用之间被 §11.5.2 与整节 §11.6 隔开；
%    (2) ch11-c 开篇「阻抗的基本定义、等效电路推导与 Nyquist／Bode 图的定义
%        写在上文 §11.5.1」这句话，在装配后指向的是**下文**，指代方向与实际相反；
%    (3) 编号 11.5.1 会出现在 11.5 之后（LaTeX 的 \subsection 计数器只按出现
%        顺序编号，本分片标题沿用分片 b 的层级，故书内编号将与源书目录不符）。
%  处置办法（三选一，均须改本分片之外的文件）：
%    (A) 在 chapters/ch11.tex 里把 ch11-d 的 \input 提到 ch11-b 之后、
%        ch11-c 之前（最小改动，编号与指代同时恢复正确）；
%    (B) 把本分片内容并入 ch11-b 末尾（需改 ch11-b，与并发作业冲突）；
%    (C) 接受现状，把本分片标题改为「阻抗谱基本原理（补）」并把 ch11-c 的
%        指代句改为「见前文」。
%  【已执行 (A)】chapters/ch11.tex 现为 a→b→d→c→e，实测渲染出的节标题序列为
%  11.1 → 11.2 → 11.3 → 11.4 → §11.5.1（本分片，编号 11.5）→ §11.5（分片 c，
%  编号 11.6）→ §11.6（分片 c，编号 11.7）。注意：本分片标题仍是 \subsection，
%  故**书内**自动编号为 11.5，与源书目录的「11.5.1」不同号。这是源书节号与
%  本项目层级不一致的通例——源书「11.5.1」以下的内容在本项目里被拆成两个
%  \subsection（原理与应用），源书节号一律用 \src{}／正文「§11.5.1」标注出处，
%  不映射到 book 计数器；分片 c 的指代句与 \srcnote 已同步改为
%  「上一节（§11.5.1）」并写明该编号指源书节号。
%  本分片按任务约定**不写 \chapterstart**，只用 \subsection，以免与装配层的
%  章计数器冲突（章计数器由 ch11-a 唯一推进，故调序不影响章号）。
%
%  【诚实纪律】本分片看不到原书页面图像。公式按「源文本 OCR 上下文 + 学科
%  标准形式」重建；凡是源文本未给而由学科标准形式补出的内容，都在正文或
%  图题里显式说明；凡数值、条件与符号约定存疑处，一律标 \tocheck{} 写明存疑点。
%  本分片**没有**核对过任何页面图像，正文中不作与此相反的表述。
% ================================================================

\subsection{阻抗谱基本原理}

电化学阻抗谱（\term{EIS}{electrochemical impedance spectroscopy}{电化学阻抗谱}）给电化学系统
施加一个\textbf{频率不同、振幅很小的交流正弦电势波}，测量交流电势与电流信号的比值（即系统的
阻抗）随正弦波频率的变化，或者是阻抗的相位角 $\varphi$ 随频率的变化。由于扰动电信号是交流信号，
电化学阻抗谱也称\textbf{交流阻抗谱}（alternating current impedance spectroscopy）。从装置上看得
更直观：用波形发生器产生一个小幅正弦电势信号，经恒电位仪施加到电化学系统上，把输出的电流或电势
信号经转换后送入锁相放大器或频谱分析仪，输出阻抗及其模量或相位角；改变正弦波的频率，就得到一
系列不同频率下的阻抗、阻抗模量与相位角\src{S2 §11.5.1，书内 p.855—856 / PDF p.343—344}
（原书图 11-32——电化学阻抗法实验装置示意图）。

\begin{probout}
阻抗实验回答的是「这套电极过程一共有几个可分辨的时间尺度、各是多大」。已知的是一条稳态极化
曲线只能给出所有步骤过电位的\textbf{总和}；要问的是能否把总响应按时间尺度拆开，分别读出溶液
电阻、双电层电容、电荷转移电阻与扩散阻抗。可观测量是每个频率下的两个实数——阻抗的实部与虚部
（或等价的模与相位角），而待定的等效电路元件数通常多于两个，因此本方法的核心风险不是「测不到」，
而是\textbf{解不唯一}：能拟合数据的电路不止一个，而只有一个是真实机理。（超）微电极把未补偿
电阻与时间常数压到极小，因而特别适合在低电导率介质中研究电化学过程，并能在很小的
$R_{\mathrm{u}}C_{\mathrm{d}}$ 时间常数下测量快速反应\src{S2 §11.4.3，书内 p.855 / PDF p.343}。
\end{probout}

\subsubsection{交流电、相量与简单元件的响应}

一个纯正弦电压可表示为

\begin{equation}
e=E\sin\omega t
\label{eq:eis-d-sine}
\end{equation}

式中 $\omega$ 为角频率，它是 $2\pi$ 乘以用 Hz 表示的常规频率值；$E$ 为幅值
\src{式（11-115），书内 p.856 / PDF p.344}。可以把电压用旋转矢量（或\textbf{相量}）表示：
观察到的电压 $e$ 是任意时刻投影在某一特定轴（通常取 $0$°）上相量的分量（见原书图 11-33）。
两个正弦信号——电流 $i$ 与电压 $e$——可以表示成以同样频率旋转的独立相量 $\dot I$ 和 $\dot E$，
它们通常不同相，相差一个相角 $\varphi$，以 $\dot E$ 为参考信号，$\varphi$ 即相对它测出。图 11-34
中电流滞后于电压，此时 $\varphi<0$，可表示为

\begin{equation}
i=I\sin(\omega t+\varphi)
\label{eq:eis-d-phase}
\end{equation}

\src{式（11-116），书内 p.857 / PDF p.345；原书图 11-34 给出频率为 $\omega$ 的交流电压与电流
相互关系的相量图}。\tocheck{未对照原页面图像：源文本此式即印作
$i=I\sin(\omega t+\varphi)$，本稿照印、不再改写成 $\sin(\omega t-\varphi)$——按本节的测量相角
约定 $\varphi\equiv-\arg Z$（见下文式~\eqref{eq:eis-d-zmod} 后的约定行），该符号是自洽的：电流
滞后即 $\varphi<0$；纯电容时式（11-118）的 $i=(E/X_C)\sin(\omega t+\pi/2)$ 给出
$\varphi=+\pi/2$。该式的符号未能逐字与页面核对。}

简单元件对电压的响应是本节的积木。纯电阻 $R$ 上施加 $e=E\sin\omega t$ 时，由欧姆定律电流为
$(E/R)\sin\omega t$，相角为零（原书图 11-35）。把电阻换成纯电容 $C$，所需的基本关系式是
$q=it$、$q=Ce$ 与 $i=C(\mathrm{d}e/\mathrm{d}t)$，因此

\begin{equation}
i=\omega CE\cos(\omega t),\qquad
i=\frac{E}{X_C}\sin\!\left(\omega t+\frac{\pi}{2}\right)
\label{eq:eis-d-cap}
\end{equation}

式中 $X_C$ 为\textbf{容抗}，其值为 $1/\omega C$\src{式（11-117）、（11-118），书内 p.857—858 /
PDF p.345—346}。相关电流的公式与电阻相似，只是以 $X_C$ 代替 $R$，且相角由 $0$ 变为 $\pi/2$：
\textbf{电流导前于电压}（原书图 11-36）。由于相角改变后矢量图扩展为一个平面，用复数符号表示
相量就成为必要：规定横坐标分量为实部、纵坐标分量为虚部，并乘以 $j=\sqrt{-1}$。这里引入复数符号，
是想使矢量的各分量都成为直线——实部与虚部这两种形式在相角可测量的意义上讲都是真实的。

\begin{equation}
\dot E=-jX_C\dot I
\label{eq:eis-d-capohm}
\end{equation}

\src{式（11-119），书内 p.858 / PDF p.346}。比较 $\dot E=\dot I R$ 与式~\eqref{eq:eis-d-capohm}
可知，$X_C$ 具有电阻的量纲，但与 $R$ 不同——\textbf{其值随频率的增加而下降}。这一点是后面所有
读图规则的根源：电容元件在频率轴上有一个确定的位置，改变频率等于改变它在电路中的「权重」。

\paragraph{串联叠加与阻抗的定义} 在电阻 $R$ 与电容 $C$ 的串联上施加电压 $\dot E$，其值无论何时
都等于电阻与电容上电压降之和：

\begin{equation}
\dot E=\dot E_R+\dot E_C=\dot I\,(R-jX_C)=\dot I Z
\label{eq:eis-d-series}
\end{equation}

电压与电流通过一个称为\textbf{阻抗}的矢量 $Z=R-jX_C$ 联系在一起
\src{式（11-120）至（11-122），书内 p.858 / PDF p.346}。更一般地，阻抗写成实部与虚部之和：

\begin{equation}
Z(\omega)=Z_{\mathrm{Re}}+jZ_{\mathrm{Im}},\qquad
Z_{\mathrm{Re}}=R,\qquad Z_{\mathrm{Im}}=-X_C=-\frac{1}{\omega C}
\label{eq:eis-d-zdef}
\end{equation}

$Z$ 的幅值与相角为

\begin{equation}
\lvert Z\rvert^{2}=R^{2}+X_C^{2}
=\left(Z_{\mathrm{Re}}\right)^{2}+\left(Z_{\mathrm{Im}}\right)^{2},
\qquad
\tan\varphi=-\frac{Z_{\mathrm{Im}}}{Z_{\mathrm{Re}}}=\frac{X_C}{R}
=\frac{1}{\omega RC}
\label{eq:eis-d-zmod}
\end{equation}

\src{式（11-123）至（11-125），书内 p.858—859 / PDF p.346—347}。式~\eqref{eq:eis-d-zmod} 中的
$\varphi$ 是\textbf{测量相角}，本节统一取 $\varphi\equiv-\arg Z$（$\arg Z$ 为 $Z$ 的辐角），
即\textbf{容性为正}：纯电阻 $\varphi=0$、纯电容 $\varphi=+\pi/2$、纯电感 $\varphi=-\pi/2$；
对混合体则观察到介于相应两者之间的相角。\textbf{Nyquist 图纵轴取 $-Z''$（向上为正）与 Bode 图的
相角曲线同属这一约定}——同一条谱在两种图上必须给出同一个 $\varphi$，换图不换号。
\textbf{更为复杂的电路，可以按与电阻相同的规则通过合并阻抗来分析——串联阻抗的总阻抗是各阻抗值
之和；并联则用导纳 $Y\equiv1/Z$ 表述：总导纳是各支路导纳之和，即}
\begin{equation}
Y_{\text{并}}=\frac{1}{Z_{\text{并}}}=\sum_i\frac{1}{Z_i},
\qquad
Z_{\text{并}}=\Bigl(\sum_i\frac{1}{Z_i}\Bigr)^{-1}
\label{eq:eis-d-parallel}
\end{equation}
\src{书内 p.859 / PDF p.347。原文作「并联阻抗的总阻抗是单个阻抗值倒数之和」，本式即按「总阻抗的
倒数是各阻抗倒数之和」理解}。核对两个数：两个 $10$ \unit{\ohm} 的电阻并联，
$Y_{\text{并}}=0.1+0.1=0.2$ S，故 $Z_{\text{并}}=1/0.2=5$ \unit{\ohm}（若把「倒数之和」
$0.2$ S 直接当成阻抗值，会得到 $0.2$ \unit{\ohm}，与 $5$ \unit{\ohm} 差 $25$ 倍）；$R$ 与
$C$ 并联时 $Y=1/R+j\omega C$，低频极限 $Z\to R$（$\varphi\to0$）、高频极限
$Z\to-j/(\omega C)$（$\varphi\to+\pi/2$）。这两条合并规则加上式~\eqref{eq:eis-d-series}，就是
全部等效电路的数学基础；后文所有电路的阻抗表达式都只是这两条规则的反复应用。
\tocheck{未对照原页面图像：源文本式（11-123）以下的 $Z_{\mathrm{Im}}$ 印作 $+X_C$、式（11-125）
印作 $\tan\varphi=X_C/R$，二者与式（11-121）$Z=R-jX_C$ 相差一个符号；本稿取
$Z_{\mathrm{Im}}=-X_C$（即 $j$ 的系数）并令 $\varphi\equiv-\arg Z$，两式即同时成立，也与源文本
p.859「纯电容 $\varphi=\pi/2$」一致。源文本式（11-125）OCR 只剩
$\tan\varphi=1/\omega RC$ 的骨架，符号与绝对值未能逐字核对。}

\begin{principle}
\textbf{把电化学系统看成等效电路。} 利用电化学阻抗谱研究一个电化学系统时，它的基本思路是把
电化学系统看作一个\textbf{等效电路}——由电阻（$R$）、电容（$C$）和电感（$L$）等基本元件按串联
或并联等不同方式组合而成；通过 EIS 测定等效电路的构成及各元件的大小，再利用这些元件的电化学
含义分析电化学系统的结构和电极过程的性质\src{书内 p.859 / PDF p.347}（原书图 11-37——电阻、
电容及电感等元件的串联和并联）。EIS 研究是在\textbf{可逆}电极反应过程的基础上发展起来的，
用线性元件作为等效元件，构成能给出与实测 EIS 一样谱图的等效电路\src{书内 p.860 / PDF p.348}。
「可逆」与「线性元件」这两个前提是整套方法的边界：它要求扰动足够小，使电流–电势关系可以线性化，
从而能够用传递函数（而不是非线性微分方程）描述系统。
\end{principle}

\subsubsection{两种常用画法与其定义}

电化学阻抗技术通过测定不同频率下扰动信号 $X$ 与响应信号 $Y$ 的比值，得到各频率下阻抗的实部、
虚部、模值和相位角，再把这些量绘成各种形式的曲线，就得到电化学阻抗谱
\src{书内 p.859 / PDF p.347}。常用的阻抗谱有两种。

\paragraph{Nyquist 图（能斯特图）} 以阻抗的实部 $Z_{\mathrm{Re}}$ 为横轴、虚部 $Z_{\mathrm{Im}}$
为纵轴，\textbf{图中的每个点代表不同的频率，左侧的频率高，成为高频区，右侧的频率低，成为低频
区}\src{书内 p.859 / PDF p.347}。本书从本节起统一按电化学惯例以 $-Z''$ 向上为正作图
（即纵轴画 $-Z_{\mathrm{Im}}$），以便常见的容性数据落在第一象限；这一点在读书中的数据图时必须
先确认，否则整个图会上下翻转。

\paragraph{Bode 图（波特图）} 包括\textbf{两条}曲线，它们的横坐标都是频率的对数 $\lg f$，纵坐标
一个是阻抗模值的对数 $\lg\lvert Z\rvert$，另一个是阻抗的相位角 $\varphi$
\src{书内 p.859 / PDF p.347}。两图的分工是本节的第一个实用结论：\textbf{Nyquist 图把频率信息
抹掉了}——图上每个点对应哪个频率必须另查数据表；\textbf{Bode 图把频率显式保留}，因而判断「哪个
频段被哪个过程支配」更可靠。因此\textbf{同一组数据必须两图并看}：用 Nyquist 图读几何量（截距、
直径、斜率），用 Bode 图定频段与数时间常数。

\paragraph{三类极限体系的特征谱} 源文本给出三种基本情况，可以当作读图的「字典」
\src{书内 p.859—860 / PDF p.347—348}：

\begin{itemize}
\item \textbf{纯电阻}：Nyquist 图上表现为 $Z'$ 轴（横轴）上的一个点，该点到原点的距离为电阻值的
大小；Bode 图中 $\lg\lvert Z\rvert$–$\lg f$ 为一条水平直线，相角 $\varphi$ 为 $0$°且不随测量
频率变化。
\item \textbf{纯电容体系}：Nyquist 图上表现为与纵轴（$-Z''$ 轴，按本节 $\varphi\equiv-\arg Z$ 的
约定向上为正）重合的一条直线；Bode 图中 $\lg\lvert Z\rvert$–$\lg f$ 是斜率为 $-1$ 的直线，
$\varphi$ 为 $+90$°。
\item \textbf{Warburg 阻抗} $Z_{\mathrm{W}}$：Nyquist 图上表现为相角为 $45$°的直线；Bode 图上
表现为斜率为 $-1/2$、$\varphi$ 为 $+45$°的直线。
\end{itemize}

这幅「字典」的价值在于它给出了\textbf{三条可直接反查的判据}：$|Z|$ 对频率的幂律指数 $0$、$-1$、
$-1/2$ 分别对应电阻、电容与半无限扩散。测得一段数据的局部斜率后，先问「它接近 $0$、$-1/2$ 还是
$-1$」，再去谈具体机理——这个顺序可以挡住大部分把 Warburg 当半圆、把 CPE 当电容的误判。

\src{式（11-115）至（11-125）依次见书内 p.856—859 / PDF p.344—347，原书图 11-33 至图 11-37
依次见书内 p.857—859 / PDF p.345—347。}

\subsubsection{等效元件与四类电极的特征谱}

等效电路中的常用等效元件如表~\ref{tab:eis-d-elements} 所示，系源文本表 11-1 的内容
\src{原书表 11-1，书内 p.860 / PDF p.348}。

\begin{table}[htbp]
\centering
\caption{等效电路中的常用等效元件，以及读图时各自的成立条件。第四列是本分片
按学科标准形式补出的：\textbf{它回答的是「凭什么把这个图形特征归到这个元件上」}，条件不成立时
读出的数不是「误差大一点」，而是物理归属错了。\src{据 S2 表 11-1，书内 p.860 / PDF p.348
整理。}}
\label{tab:eis-d-elements}
\begin{tabularx}{\onecolwidth}{@{}p{17mm}p{8mm}p{40mm}X@{}}
\toprule
名称 & 符号 & 源文本给出的意义 & 读图用法与成立条件 \\
\midrule
溶液电阻 & $R_{\Omega}$ & 工作电极与参比电极之间的电阻 &
$R_{\Omega}$ 与接触电阻、引线电阻在谱上不可分辨——EIS 本身分不开它们。故
「高频截距 $=R_{\Omega}$」是\textbf{约定}而非测量结果：须另用改变鲁金毛细管位置（截距随之变）或
电流中断法标定，才可把截距当溶液电阻用。 \\
双电层电容 & $C_{\mathrm{dl}}$ & 工作电极与电解质之间的电容 &
与 $R_{\mathrm{ct}}$ 并联，构成一个时间常数 $\tau=R_{\mathrm{ct}}C_{\mathrm{dl}}$。仅在电位区间内
无法拉第电流时，它才单独决定响应（此时半圆不出现，谱退化为垂直线）。 \\
极化电阻 & $R_{\mathrm{p}}$ & 当电位远离开路电位时，导致电极表面产生电流，电流受到反应动力学和
反应物扩散的控制 &
$R_{\mathrm{p}}$ 是\textbf{动力学与扩散的总和}，不是纯动力学量；只在扩散可忽略时
$R_{\mathrm{p}}\approx R_{\mathrm{ct}}$。用它算 $i_0$ 前必须先确认低频没有扩散尾。 \\
电荷转移电阻 & $R_{\mathrm{ct}}$ & 电化学反应动力学控制 &
对应半圆直径，但须「弧闭合、只有一个主时间常数、扩散不参与」；三者缺一，直径里就混着别的贡献。 \\
扩散电阻（Warburg 阻抗） & $Z_{\mathrm{W}}$ & 反应物从溶液本体扩散到电极反应界面的电阻 &
对应低频 $45$°直线；成立条件是\textbf{半无限}线性扩散——扩散层厚度远小于体系最小边界尺度，且
低频点没有因体系漂移而失真。 \\
界面电容 & $C_{\mathrm{c}}$ & 通常每一个界面之间都会存在的一个电容 &
与 $C_{\mathrm{dl}}$ 并联出现（涂层、晶界、催化剂层／膜界面）。二者时间常数接近时只能读出两者的
组合量，无法按单个界面分开。 \\
电感 & $L$ & 影响电流的变化 &
图上表现为第四象限的弧。仪器电感应与频率无关、且随接线方式改变——能靠改接线把它改掉的，就不是
化学信息。 \\
\bottomrule
\end{tabularx}

\noindent{\srcsize 表中「意义」一列逐条取自源文本表 11-1（书内 p.860 / PDF p.348）；第四列为
按学科标准形式补出的操作判据，并已把源文本未给出的分界条件写明——须与原书 p.860—861 的四类
电极讨论对照。}
\end{table}

\paragraph{四类电极的 Nyquist 图} 理想极化电极、溶液电阻可以忽略的电化学极化电极、溶液电阻不能
忽略的电化学极化电极、以及电化学极化和浓差极化同时存在的电极，其阻抗谱均可用 Nyquist 图与 Bode 图
表示；源文本以 Nyquist 图为代表逐一介绍\src{书内 p.860 / PDF p.348}。

\begin{enumerate}
\item \textbf{理想极化电极}是不发生电极反应的电极，其 Nyquist 图为一条距 $Z''$ 轴为 $R_{\Omega}$
且垂直于实轴 $Z'$ 的直线；由直线在 $Z'$ 轴上的交点到原点的距离可求得 $R_{\Omega}$（原书图 11-38）
\src{书内 p.860 / PDF p.348}。
\item \textbf{溶液电阻可以忽略的电化学极化电极}对应一半径为 $R_{\mathrm{ct}}/2$ 的半圆，
由圆的直径可求得 $R_{\mathrm{ct}}$；\textbf{圆的顶点对应的 $Z''$ 最大，由对应的角频率 $\omega$
及得到的反应电阻，可求得双电层电容}
\begin{equation}
C_{\mathrm{dl}}=\frac{1}{\omega_{\max}R_{\mathrm{ct}}},
\qquad
\omega_{\max}=2\pi f_{\max}=\frac{1}{R_{\mathrm{ct}}C_{\mathrm{dl}}}
\label{eq:eis-d-cdl}
\end{equation}
（原书图 11-39）\src{书内 p.860—861 / PDF p.348—349}。\tocheck{未对照原页面图像：源文本此句
OCR 呈「由对应的角频率 $\omega$ 及得到的反应电阻，可以求得双电层电容 $C_{\mathrm{dl}}=1/(\omega
R_{\mathrm{ct}})$」，未写明是取半圆顶点的 $\omega$；本分片按顶点条件补出 $\omega_{\max}$ 的
定义，并保留原式形式。（若误把 $f_{\max}$ 当作 $\omega_{\max}$ 代入，结果会差 $2\pi$ 倍——
这是本节最容易犯的算术错误。）}
\item \textbf{溶液电阻不能忽略的电化学极化电极}：设扩散过程引起的阻抗可以忽略、电极过程由电荷
传递过程（电化学反应步骤）控制，所得 Nyquist 图同样是一半径为 $R_{\mathrm{ct}}/2$ 的半圆，
由其可求得溶液电阻 $R_{\Omega}$ 和电荷转移电阻 $R_{\mathrm{ct}}$，由半圆顶点的 $\omega$ 可求得
$C_{\mathrm{dl}}$（原书图 11-40）\src{书内 p.861 / PDF p.349}。
\item \textbf{电化学极化和浓差极化同时存在的电极}：电极过程由电荷传递过程和扩散过程共同控制，
其 Nyquist 图由\textbf{高频区一个半径为 $R_{\mathrm{ct}}/2$ 的半圆}和\textbf{低频区一条 $45$°的
直线}构成——半圆表示电荷传递过程为控制步骤，直线表示由电极反应的反应物或产物的扩散控制
（原书图 11-41）\src{书内 p.861 / PDF p.349}。
\end{enumerate}

这四类讨论给出的正是三条最常用的读数规则。但源文本是按「理想情形」逐类陈述的，读者若把每一类
的结论直接搬到实测数据上，就会丢掉前提。下面把三条规则各自的成立条件写全，它们分别对应
「高频截距」「半圆直径」「低频斜线」三个图形特征。

\paragraph{规则一：高频截距 $=R_{\Omega}$} 成立需要同时满足两条：
（1）\textbf{频率足够高}，高到 $C_{\mathrm{dl}}$ 的阻抗小到可以忽略、$R_{\mathrm{ct}}$ 被完全旁路，
即 $\omega R_{\mathrm{ct}}C_{\mathrm{dl}}\gg1$；
（2）\textbf{高频极限处的接触电阻与引线电感可以独立排除}——接触电阻在物理上与溶液电阻串联，
EIS 无法从谱上把它与 $R_{\Omega}$ 分开，只能靠改变电极压紧力或参比位置、看截距是否跟着变来判定；
引线电感则使实部先减小后增大，把截距抬偏甚至读不出。两条都不满足时，截距里含 $R_{\mathrm{ct}}$
成分，读出的「溶液电阻」偏大。

\paragraph{规则二：半圆直径 $=R_{\mathrm{ct}}$} 成立需要同时满足三条：
（1）\textbf{半圆闭合}，即频率范围覆盖了从 $\omega\tau\gg1$ 到 $\omega\tau\ll1$ 的全过程；
（2）\textbf{这段弧只由一个时间常数支配}——两个时间常数相差不到约两个数量级时，两个半圆会重
叠成一个压扁的弧，此时用「直径」读出的既不是 $R_{\mathrm{ct}1}$ 也不是 $R_{\mathrm{ct}2}$；
（3）\textbf{扩散不参与}，否则弧的低频一侧被 Warburg 尾「拉住」，外推得到的截距偏大。
要注意\textbf{压扁也可能来自 CPE 而非两个过程}，二者仅凭一张 Nyquist 图无法区分，必须看 Bode
相角图上有几个峰。

\paragraph{规则三：低频 $45$°斜线 $=$ 扩散控制} 成立需要在\textbf{两端}同时成立：
（1）\textbf{上界}——频率低到 Warburg 项主导，即 $\sigma\omega^{-1/2}$ 已远大于 $R_{\mathrm{ct}}$
的贡献；
（2）\textbf{下界}——频率还不够低，扩散层厚度 $\sqrt{2D/\omega}$ 仍\textbf{远小于}体系的最小
边界尺度（薄液层厚度、膜厚、多孔层厚度、对流边界层厚度）。把这两个条件写成可算的判据：半无限
扩散的有效下限约在
\begin{equation}
f_{\mathrm{low}}\gtrsim\frac{D}{2\pi\delta^{2}}
\label{eq:eis-d-fdiff}
\end{equation}
$f<f_{\mathrm{low}}$ 之后扩散层「看到」了边界，$45$°直线必然弯曲（弯向实轴或转为接近竖直的
电容性直线）。\tocheck{未对照原页面图像：式~\eqref{eq:eis-d-fdiff} 是本分片按半无限扩散的
特征扩散长度 $\delta\sim\sqrt{2D/\omega}$ 补出的操作性判据，源文本 §11.5.1 未给出该式；其常数
因子（$2\pi$ 与 $2$ 的取法）取决于所用特征长度的定义，须与 §11.5.1 的符号约定对照。}
数量级感受：$D=1\times10^{-5}$ cm$^{2}$·s$^{-1}$ 且薄液层厚 $\delta=100$ µm 时
$f_{\mathrm{low}}\approx0.03$ Hz——也就是\textbf{常规阻抗实验从 10 mHz 起扫
的那些低频点，早已越过了薄层体系的半无限下限}。在不搅拌的厚层溶液中（边界是本体溶液）这一条
通常自动满足；在薄膜电极、薄液层电解池与多孔催化剂层里，它是必须先算一遍的条件。

\subsubsection{等效电路的阻抗表达式（Randles 电路、Warburg 阻抗与 CPE）}

源文本把等效电路作为「基本思路」提出，并给出表 11-1 的元件清单与四类电极的 Nyquist 图，但
\textbf{没有写出这些电路的总阻抗表达式}。下面把最常用的几个表达式按学科标准形式补出，作为读图
（下一小节）的工具；凡源文本已给出形式的，都在公式后标出页码，凡由标准形式补出的，逐条说明。

\paragraph{Randles 电路与溶液电阻、电荷转移电阻、双电层电容} 把溶液电阻 $R_{\Omega}$（源文本符号，
即通常写作 $R_{\mathrm{s}}$ 者）与一个并联支路串联，并联支路由电荷转移电阻 $R_{\mathrm{ct}}$ 与
扩散阻抗 $Z_{\mathrm{W}}$ 串联后，再与双电层电容 $C_{\mathrm{dl}}$ 并联，即得最常用的
\textbf{Randles 电路}：

\begin{equation}
Z=R_{\Omega}+\left[\frac{1}{R_{\mathrm{ct}}+Z_{\mathrm{W}}}
+j\omega C_{\mathrm{dl}}\right]^{-1}
\label{eq:eis-d-randles}
\end{equation}

\tocheck{未对照原页面图像：Randles 电路的具体拓扑（$Z_{\mathrm{W}}$ 与 $R_{\mathrm{ct}}$ 串联、
二者整体与 $C_{\mathrm{dl}}$ 并联）与式~\eqref{eq:eis-d-randles} 是按源文本图 11-38 至图 11-41
的四类电极 Nyquist 图所隐含的结构、以及学科标准形式重建的；源文本 §11.5.1 未写出该式，故拓扑
本身未能与页面核对。}
当扩散可忽略（$Z_{\mathrm{W}}\to0$）时，式~\eqref{eq:eis-d-randles} 退化为
$R_{\Omega}$ 串联一个并联 $R_{\mathrm{ct}}C_{\mathrm{dl}}$ 网络，即规则二讨论的单个半圆；此时
\begin{equation}
Z'=R_{\Omega}+\frac{R_{\mathrm{ct}}}{1+(\omega R_{\mathrm{ct}}C_{\mathrm{dl}})^{2}},
\qquad
-Z''=\frac{\omega R_{\mathrm{ct}}^{2}C_{\mathrm{dl}}}{1+(\omega R_{\mathrm{ct}}C_{\mathrm{dl}})^{2}}
\label{eq:eis-d-rc}
\end{equation}
两式消去 $\omega$ 即得圆心在实轴上 $R_{\Omega}+R_{\mathrm{ct}}/2$ 处、半径为 $R_{\mathrm{ct}}/2$
的圆方程——这是「半圆直径 $=R_{\mathrm{ct}}$」与「半圆顶点频率
$\omega_{\max}=1/(R_{\mathrm{ct}}C_{\mathrm{dl}})$」两条读图规则的来源。\textbf{两条规则的几何
前提是同一个}：谱必须是一个完整的半圆；谱被压扁时式~\eqref{eq:eis-d-rc} 不再成立，圆的圆心
落到实轴下方，直径也不再等于 $R_{\mathrm{ct}}$。

\paragraph{Warburg 阻抗与扩散系数} 半无限线性扩散给出

\begin{equation}
Z_{\mathrm{W}}=\sigma\,\omega^{-1/2}(1-j),
\qquad
Z'_{\mathrm{W}}=-Z''_{\mathrm{W}}=\sigma\,\omega^{-1/2},
\qquad
\sigma=\frac{RT}{n^{2}F^{2}A\sqrt{2}}
\left(\frac{1}{c_{\mathrm{O}}^{*}\sqrt{D_{\mathrm{O}}}}
+\frac{1}{c_{\mathrm{R}}^{*}\sqrt{D_{\mathrm{R}}}}\right)
\label{eq:eis-d-warburg}
\end{equation}

实部与虚部恒相等，故相角恒为 $+45$°（按 $\varphi\equiv-\arg Z$），在 Nyquist 图上是一条斜率 $1$ 的直线
\src{相角 $45$°这一结论见书内 p.859 / PDF p.347}。由 $45$°直线在 $\omega^{-1/2}$ 轴上的斜率读出
$\sigma$ 后，若已知 $n$、$A$ 与浓度，即可估出扩散系数 $D$——这是阻抗谱少有的「能给出传质定量
参数」的场合。\tocheck{未对照原页面图像：$\sigma$ 的表达式（含系数 $RT/(n^{2}F^{2}A\sqrt{2})$
与括号内两项）以及 $Z_{\mathrm{W}}$ 的符号约定按学科标准形式重建，源文本 §11.5.1 只给出
「相角 $45$°」这一图形结论，未给出 $\sigma$ 的定义式，故未能与页面核对。}

\paragraph{常相位元件（CPE）} 实测半圆几乎从不以实轴为直径，而是一个\textbf{压扁的弧}。把
$C_{\mathrm{dl}}$ 换成\textbf{常相位元件}（constant phase element，CPE）即可描述这一现象：

\begin{equation}
Z_{\mathrm{CPE}}=\frac{1}{Q(j\omega)^{n}},\qquad 0<n\le1,
\qquad
\varphi_{\mathrm{CPE}}=\frac{n\pi}{2}
\label{eq:eis-d-cpe}
\end{equation}

$n=1$ 时 $Q$ 就是电容（单位为 F）；$n=0$ 时退化为电阻。\textbf{$n$ 偏离 $1$ 的物理含义是电极
表面不均匀}：粗糙表面、多晶面的不同晶面取向、吸附层覆盖不均、以及电流分布不均，都会使界面的
弛豫时间不是一个值，而是一段分布——分布越宽，$n$ 越小、弧越扁。\textbf{把 CPE 当理想电容去算
电容值，是在这一节里最容易发生的系统性错误}：$Q$ 的量纲是 $\mathrm{S.s^{n}}$，它不是电容；由
$Q$ 与 $n$ 反求有效电容必须做量纲换算
\begin{equation}
C_{\mathrm{eff}}=\left(Q\,R_{\mathrm{ct}}^{\,1-n}\right)^{1/n}
\label{eq:eis-d-cpeceff}
\end{equation}
而有效电容的取法（是否用 $R_{\Omega}$ 还是 $R_{\mathrm{ct}}$、或两者的并联）本身有不同约定，
因此\textbf{报告结果时必须同时给出 $n$ 与所用换算式的出处}，否则读者无法复算。
\tocheck{未对照原页面图像：式~\eqref{eq:eis-d-cpe} 中 $Q$ 的符号与 $0<n\le1$ 的区间、以及
式~\eqref{eq:eis-d-cpeceff} 的换算形式按学科标准形式重建，源文本 §11.5.1 未出现 CPE 或 $n$ 的
讨论（原书表 11-1 只有理想双电层电容 $C_{\mathrm{dl}}$），故该元件与原书的对照关系未能核对。}

\begin{principle}
\textbf{为什么等效电路能成立，以及它成立的代价。} 式~\eqref{eq:eis-d-randles} 一类的表达式之所以
可以描述一个真实的电极／溶液界面，靠的是三条前提：\textbf{线性}（扰动足够小，$i$–$E$ 关系可在
工作点附近线性化）、\textbf{因果}（响应不先于扰动）、\textbf{稳定}（测量期间体系不随时间漂移，
即处于稳态）。三条合起来等价于「系统可用线性时不变传递函数描述」，这正是阻抗定义
$Z(\omega)=E(\omega)/I(\omega)$ 的适用条件。Kramers--Kronig 关系是这三条前提的数学后果，因此
它可以反过来被用作\textbf{数据合格性的检验}：残差在全部频率上都小，才说明上面三条成立；否则
该组数据不应进入拟合\src{书内 p.859 / PDF p.347 给出「观察体系在稳态时对扰动的响应，是频率域
的测量」这一前提；Kramers--Kronig 检验为其数学后果，属本分片按学科标准形式补充}。
代价也要说清：\textbf{等效电路是等价描述，不是唯一描述}。同一组数据可以有多个电路满足它，
「线性元件」也只是在工作点附近成立——因此阻抗谱给出的是\textbf{同一稳态下}各步骤的时间尺度与
权重，它不能单独回答各步骤之间的先后关系。
\end{principle}

\subsubsection{测量装置与测量条件}

\begin{measure}
\textbf{（1）装置链路。} 波形发生器产生小幅正弦电势 $\to$ 恒电位仪施加到电化学系统 $\to$ 输出的
电流／电势信号经转换 $\to$ 锁相放大器或频谱分析仪 $\to$ 输出阻抗及其模量或相位角；改变正弦波
频率即得一系列不同频率下的阻抗、模量与相位角\src{书内 p.855—856 / PDF p.343—344}。链路上每一环
都会在某个频段引入自己的误差：恒电位仪的带宽与相位误差在高频端，锁相放大器的时间常数在低频端，
电缆与接线的电感与电容在高频端。\textbf{因此每次实验都应先用已知的 dummy cell（纯电阻或
$R$–$C$ 串联）验证整条链路}，特别是相位——相位误差不能靠多测平均消除。

\textbf{（2）扰动幅值。} 必须小到使响应线性，又大到使信噪比足够。判据不是记住一个数字，而是做
\textbf{幅度减半检验}：把交流幅值减半后重测，若全频段谱形不变，说明仍在线性区；若 $R_{\mathrm{ct}}$
随幅值系统性变小，说明已进入非线性区。燃料电池测量的交流振幅一般选择在 $5\sim15$ mV
\src{书内 p.862 / PDF p.350}；催化剂活性测量因电位控制更直接，通常取更小值。

\textbf{（3）频率范围。} 范围本身是结论的一部分。燃料电池测量中扫描频率范围一般在
1 mHz～100 kHz 之间\src{书内 p.862 / PDF p.350}。范围的选择须由
\textbf{待求的时间常数}反推：高频端要跨过 $1/(2\pi R_{\mathrm{ct}}C_{\mathrm{dl}})$ 至少两个
数量级（否则高频截距含 $R_{\mathrm{ct}}$ 成分），低频端要覆盖目标过程（扩散尾要看到 $45$°、
吸附弛豫要看到相角峰）。按照规则三下界的量级估算，低频端的点数应密集到足以分辨斜率，而不是按
对数均匀取点。

\textbf{（4）测量模式。} 阻抗测量一般有\textbf{恒电位模式}（扰动为交流电位，响应为测量电流）与
\textbf{恒电流模式}（扰动为交流电流，响应为测量电位）两种，从两种模式得到的结果基本上没有明显
不同，但各有优缺点\src{书内 p.862 / PDF p.350}。在恒电位模式中，一个小的电位都可能导致大的电流
变化，从而导致恒电位仪过载，因此尝试不同的交流电压以确保恒电位仪不过载是必要的；在催化剂活性
的测量中，控制电位能更好地控制电化学反应的发生，因此一般采用恒电位模式。在恒电流模式中，交流
和直流大小很容易被运算放大器确定，能有效防止电池和恒电位仪过载，因此能采用相对大的交流电，
从而得到更多的阻抗信息；对于电池阻抗测量，一般电流比较大，更适合用恒电流模式，恒电位模式容易
导致过载\src{书内 p.862 / PDF p.350}。

\textbf{（5）稳态必须先行建立。} 阻抗定义要求体系在测量期间是线性时不变系统，因此
\textbf{「稳态未建立就测」是本节一切误判的共同源头}。必须先在目标电位恒电位极化到电流漂移小于
约每分钟 $1\%$，再开始扫频；对慢过程（覆盖度还在演化、催化剂正在失活），应先在\textbf{一个
固定频率上监测 $\lvert Z\rvert$ 随时间的变化}，待其稳定后再扫全谱。\tocheck{未对照原页面图像：
「电流漂移小于约每分钟 $1\%$」与「先在单一频率监测 $\lvert Z\rvert$」两条操作判据按学科标准
形式补出，源文本 §11.5.1 只给出「观察体系在稳态时对扰动的响应」这一前提，未给出量化判据。}

\textbf{（6）记录条件。} 报告阻抗结果必须同时给出：交流幅值与模式、频率范围与每十倍频程的点数、
直流偏置电位、电极面积与几何、溶液组成与温度、参比电极种类、以及是否扣除与如何扣除
$R_{\Omega}$。缺任何一项，别处无法复现，也无法判断读出的数是否可比。
\src{本节（1）（2）（4）三条中凡标页码处为源文本 §11.5.2 对测量条件的直接陈述
（书内 p.862 / PDF p.350）；本分片只取与 §11.5.1 原理直接相关的条件，应用层面的讨论
（膜透过、催化剂、MEA、单池、电堆）不在本节范围内。}
\end{measure}

\subsubsection{从一组阻抗数据到等效电路：数据流与读图实例}

\begin{dataflow}
\textbf{第一层（构造）：} 由相量关系 $Z=Z'+jZ''$ 把每个频率的一对实数量组织成一个复数；
$|Z|^{2}=(Z')^{2}+(Z'')^{2}$，$\tan\varphi=-Z''/Z'$\src{式（11-123）至（11-125），
书内 p.858—859 / PDF p.346—347}。源文本在此处把 $Z_{\mathrm{Re}}$ 记作 $R$、把虚部印作
$X_C$（符号不严，见上文式~\eqref{eq:eis-d-zmod} 后的 \tocheck），那是式（11-121）的串联 $RC$
情形；本节其余处一律用一般记号 $Z'$、$Z''$，相角取测量约定 $\varphi\equiv-\arg Z$。

\textbf{第二层（定向）：} Nyquist 图上横轴 $Z'$ 向右增大、纵轴 $-Z''$ 向上增大，
\textbf{左侧频率高、右侧频率低}\src{书内 p.859 / PDF p.347}。因此沿曲线从左端走到右端 =
频率从高降到低；这条约定是后面所有「先看到谁」判断的基础。

\textbf{第三层（不变量检验）：} 纯电阻对应 $|Z|$–$f$ 水平线与 $\varphi=0$°；纯电容对应斜率
$-1$ 与 $\varphi=+90$°；Warburg 阻抗对应斜率 $-1/2$ 与 $\varphi=+45$°
\src{书内 p.859—860 / PDF p.347—348}。\textbf{先在 Bode 图上量出局部幂律指数，再回 Nyquist
图谈几何}——顺序颠倒就会把 Warburg 的直线段当圆弧读。

\textbf{完整数据流：} 逐频率数据 $\to$（检验）Kramers--Kronig 残差与正反扫重合性 $\to$
（变换）$\lg\lvert Z\rvert$ 与 $\varphi$ 对 $\lg f$ 作图、数相角峰个数定时间常数个数 $\to$
（分频段）高频平台定 $R_{\Omega}$、中频弧定 $R_{\mathrm{ct}}$ 与 $C_{\mathrm{dl}}$（或 $Q$ 与
$n$）、低频尾定扩散 $\to$（拟合）选定电路拓扑做非线性最小二乘 $\to$（验收）残差随机分布、
参数置信区间合理、跨电位／跨转速趋势自洽 $\to$（归属）把每个参数接到一个物理量与一条成立
条件上。\textbf{这条链上任何一步被跳过，后面的结论都不成立}；其中「数相角峰个数」是唯一能定
「电路里至少需要几个时间常数」的步骤，不能由拟合替代。
\end{dataflow}

\begin{figure}[tbp]
\centering
\begin{tikzpicture}
\tikzset{declare function={
  dpw(\z)  = 10^(\z);
  dwn(\z)  = 2*3.14159265358979*dpw(\z);
  dtc(\z)  = 100*0.5e-6;
  dpr(\z)  = dwn(\z)*dtc(\z);
  didx(\z) = 10+100/(1+(dpr(\z))^2);
  didy(\z) = 100*dpr(\z)/(1+(dpr(\z))^2);
  dq0      = (0.5e-6)^(0.85)*100^(0.85-1);
  dcpn     = 0.85;
  % 角度一律用「度」：pgfmath 的 sin/cos/tan 取度不取弧度（R3-C02 的实测探针
  %   sin(0.85*pi/2)=0.02328 是 sin(1.335°) 的值）。原先写 dcpb=n*pi/2 等于把
  %   θ 当成 1.335 rad 用，sin/cos 全部算错，弧被压成一条贴着横轴的线。
  dcpb     = dcpn*90;
  % u = Rct*Q*ωⁿ 才是无量纲组合；原写法 (ω*100*Q)^n 把整个乘积升到 n 次，
  %   量纲和数值都不对。下面的 a、d 分别是分母 1+u(cosθ+j sinθ) 的实部与虚部。
  dcpu(\z) = 100*dq0*(dwn(\z))^(dcpn);
  dcpa(\z) = 1+dcpu(\z)*cos(dcpb);
  dcpd(\z) = dcpu(\z)*sin(dcpb);
  dcpx(\z) = 10+100*dcpa(\z)/((dcpa(\z))^2+(dcpd(\z))^2);
  dcpy(\z) = 100*dcpd(\z)/((dcpa(\z))^2+(dcpd(\z))^2);
  % CPE 弧顶的解析值：对 u 求极值（u=1）得 Z'=Rs+Rct/2、-Z''=(Rct/2)tan(nπ/4)。
  %   标记点由这两个式子给出，不再手写孤立坐标。
  dcpxp    = 60;
  dcpyp    = 50*tan(dcpn*45);
  % Randles 加半无限 Warburg：Z = Rs + [1/(Rct+σ(1-j)/√ω) + jωC]^{-1}
  %   （式~\eqref{eq:eis-d-randles}）。原写法把它当成 Rct∥Warburg 的两元并联
  %   —— 丢失了 Cdl 支路 —— 并且把有量纲的 σ/√ω 直接与 1 相加（R3-C03 的根因）。
  %   下面改用归一化导纳 Ŷ = Y·Rct（无量纲）：
  %     t = σ/(Rct√ω)，S = (1+t)²+t²，Ŷ = (1+t)/S + j[t/S + ω C Rct]
  %   再取 Z - Rs = Rct/Ŷ。中间量都远离 pgfmath 的数值上下限，低频段也不会
  %   因为两个小量相减而丢精度。
  dkwz(\z) = 40/(100*sqrt(dwn(\z)));
  dkwS(\z) = (1+dkwz(\z))^2+(dkwz(\z))^2;
  dkwA(\z) = (1+dkwz(\z))/dkwS(\z);
  dkwB(\z) = dkwz(\z)/dkwS(\z)+dwn(\z)*0.5e-6*100;
  dkwN(\z) = (dkwA(\z))^2+(dkwB(\z))^2;
  dkwx(\z) = 10+100*dkwA(\z)/dkwN(\z);
  dkwy(\z) = 100*dkwB(\z)/dkwN(\z);
  % R2-C06：原先还声明过 dm / dmt / dmph / diph 四个量，已删除——
  %   ① 全书 grep 确认它们从未被任何 \addplot 调用（纯死代码）；
  %   ② 数学上也不成立：`diph` 给出的是 atan(ωRC)，按本节 φ≡−arg Z
  %      （tanφ=1/ωRC）恰是 φ 的余角；并且 pgfmath 的三角函数取「度」，
  %      要算 tan(76.5°) 就得写 tan(76.5)，写成 tan(dm*90) 之类是弧度思维。
  %   留着会被后来者当成可用函数，故删除而不是改写。本图不画相角曲线。
}}

\begin{axis}[
  name=eisda,
  width=70mm, height=62mm,
  unit vector ratio=1 1 1,
  xlabel={$Z'$ / $\Omega$}, ylabel={$-Z''$ / $\Omega$},
  % 轴伸到 62 Ω（刻度仍到 60），让两个弧顶标记落在框内而不是压在边框上。
  xmin=0, xmax=62, ymin=0, ymax=62,
  xtick={0,20,40,60}, ytick={0,20,40,60},
]
  \addplot[mainline,black,domain=0.3:4.9,samples=220,variable=\z]
    ({didx(\z)},{didy(\z)});
  \addplot[mainline,black,dashed,domain=0.3:4.9,samples=220,variable=\z]
    ({dcpx(\z)},{dcpy(\z)});
  \addplot[only marks,mark=*,mark size=1pt,black]
    coordinates {(60,50)};
  \addplot[only marks,mark=square*,mark size=1pt,black]
    coordinates {({dcpxp},{dcpyp})};
  \node[anchor=north west,font={\fontsize{7}{8.5}\selectfont}]
    at (axis cs:1.2,60.5) {(a)};
  % 曲线身份直接写在面板里（理想弧在 y>50 Ω 处、CPE 弧在 x>22 Ω 处都已是最上沿，
  %   这两处标注不会与任何曲线交叉）。图例因此取消，避免文字块压住弧顶。
  \node[anchor=west,font={\fontsize{7}{8.5}\selectfont}] at (axis cs:14,58) {实线 $n{=}1$};
  \node[anchor=west,font={\fontsize{7}{8.5}\selectfont}] at (axis cs:14,52) {虚线 $n{=}0.85$};
  % 说明块放进右下角的空白区（两条弧在 x>22 Ω、y<15 Ω 处都没有线），
  %   不再放到坐标区外，避免被裁掉。
  \node[anchor=south east,align=right,text width=28mm,font={\fontsize{7}{8.5}\selectfont}]
    at (axis cs:60.5,14.5)
    {理想弧顶点 $(60,\,50)$\\[1pt]
     CPE 弧顶点 $(60,\,39.4)$\\[1pt]
     $=(R_{\mathrm{ct}}/2)\tan(n\pi/4)$};
\end{axis}

\begin{axis}[
  name=eisdb,
  at={(eisda.right of east)}, anchor=left of west, xshift=6mm,
  width=74mm, height=60mm,
  xlabel={$Z'$ / $\Omega$}, ylabel={$-Z''$ / $\Omega$},
  xmin=0, xmax=300, ymin=0, ymax=260,
  xtick={0,50,100,150,200,250,300}, ytick={0,50,100,150,200,250},
]
  \addplot[mainline,black,densely dotted,domain=0.3:5.4,samples=300,variable=\z]
    ({didx(\z)},{didy(\z)});
  \addplot[mainline,black,domain=-1.7:5.4,samples=320,variable=\z]
    ({dkwx(\z)},{dkwy(\z)});
  % 六个读图点的坐标同样来自上面的函数（log10 f = -1.69897、-1、1.50515、
  %   3.50515、4.60206、4.90309，即 0.02 Hz、0.1 Hz、32 Hz、3.2 kHz、40 kHz、
  %   80 kHz），不是手写读数。
  \addplot[only marks,mark=o,mark size=1.1pt,black,mark options={fill=white}]
    coordinates {({dkwx(-1.69897)},{dkwy(-1.69897)})
                 ({dkwx(-1)},{dkwy(-1)})
                 ({dkwx(1.50515)},{dkwy(1.50515)})
                 ({dkwx(3.50515)},{dkwy(3.50515)})
                 ({dkwx(4.60206)},{dkwy(4.60206)})
                 ({dkwx(4.90309)},{dkwy(4.90309)})};
  % 45° 参考线只画一小段，放在扩散尾下方、说明块左侧的空白处（原先横跨全图，
  %   从三块说明文字上穿过去 —— R3-U01）。
  \addplot[mainline,black,dashed] coordinates {(140,10) (170,40)};
  \node[anchor=north east,font={\fontsize{7}{8.5}\selectfont}]
    at (axis cs:296,256) {(b)};
  % 面板 (b) 里最高的一条曲线只到 112.8 Ω，故 y>150 Ω 的整条带都是空白。
  \node[anchor=north west,align=left,text width=30mm,font={\fontsize{7}{8.5}\selectfont}]
    at (axis cs:8,250)
    {点线：$Z_{\mathrm{W}}{=}0$\\ 实线：$\sigma{=}40$\\ 虚线：$45$° 参考线};
  % 两块说明都落在空白区：上方（y>120 Ω 无曲线）与右下方（扩散尾下方）。
  \node[anchor=north west,align=left,text width=34mm,font={\fontsize{7}{8.5}\selectfont}]
    at (axis cs:8,182)
    {实线与点线在 $f\gtrsim 100$ Hz 几乎重合：只看高频段，\\ 分不开「$Z_{\mathrm{W}}$ 大」与「$R_{\mathrm{ct}}$ 大」};
  \node[anchor=south east,align=right,text width=22mm,font={\fontsize{7}{8.5}\selectfont}]
    at (axis cs:296,6)
    {高频截距 $\approx R_{\Omega}$\\ 只在弧闭合、\\ 无电感时成立};
\end{axis}
\end{tikzpicture}
\caption{三类等效电路的 Nyquist 图对照（模型计算）。取
$R_{\Omega}=10$ \unit{\ohm}、$R_{\mathrm{ct}}=100$ \unit{\ohm}、
$C_{\mathrm{dl}}=0.5$ µF（$\omega_{\max}=1/(R_{\mathrm{ct}}C_{\mathrm{dl}})$
对应 $f_{\max}\approx3.2$ kHz）；CPE 取 $n=0.85$、
$Q=C_{\mathrm{dl}}^{\,n}R_{\mathrm{ct}}^{\,n-1}\approx2.21\times10^{-6}$ S·s$^{0.85}$，
即低频极限电容与理想情形相同；Warburg 取 $\sigma=40$ \unit{\ohm.s^{-1/2}}。
三条曲线（实线、虚线、点线）分别由式~\eqref{eq:eis-d-randles}、
式~\eqref{eq:eis-d-cpe}、式~\eqref{eq:eis-d-rc} 按同一组参数算出：图中两个弧顶标记与
六个空心圆点都是这些函数的取值，不是手工写上的读数。
\textbf{(a)} 只画 $f=2$ Hz 至 79 kHz、$Z'\le62$ \unit{\ohm} 的窗口（两条弧的右半段落在
窗口外）：理想弧（$n=1$）的顶点在
$(R_{\Omega}+R_{\mathrm{ct}}/2,\;R_{\mathrm{ct}}/2)=(60,\,50)$，CPE 弧的顶点在
$(60,\,39.4)$——\textbf{横坐标相同，纵坐标只有 $(R_{\mathrm{ct}}/2)\tan(n\pi/4)$；
同样一个 $R_{\mathrm{ct}}$、同样一个低频极限电容，仅因 $n=0.85$ 就把弧的高度压到
$0.79$ 倍，而顶点频率并不改变}。两条曲线在 $f\to0$ 时都收敛到
$R_{\Omega}+R_{\mathrm{ct}}=110$ \unit{\ohm}（点线在 (b) 中可见 $10\to110$ \unit{\ohm}
的完整弧；CPE 弧的右半段落在 (a) 的窗口之外）。
\textbf{(b)} 频率 0.02 Hz 至 250 kHz：实线在 $f\lesssim 10$ Hz 之后已明显
偏离点线，朝向 $45$°参考线（虚线）——但\textbf{即使在 40 kHz，实线的截距仍是
10.6 \unit{\ohm} 而非 10.0 \unit{\ohm}、相角仍有 $+36.6$°（按 $\varphi\equiv-\arg Z$），即高频平台并未形成}，
这正是规则一第（1）条的量化含义；六个空心圆点依次标出后文读图实例所用的频率
（80 kHz、40 kHz、3.2 kHz、32 Hz、0.1 Hz、0.02 Hz），可以逐点核对。
注意 (b) 两个坐标轴的刻度尺度不同（横轴 300 \unit{\ohm} 与纵轴 260 \unit{\ohm}
画在等宽的图面上），故 $45$°参考线在图上并不呈 $45$°；判斜率请用数值而非图面角度。
局限：本图不含接触电阻、引线电感、有限扩散与吸附弛豫，真实数据在高频端与最低频端都会多出
本图没有的结构。
\modelmark\pedset{$R_{\Omega}$、$R_{\mathrm{ct}}$、$C_{\mathrm{dl}}$、$n$、$\sigma$ 为教学设定，
按式~\eqref{eq:eis-d-randles}、式~\eqref{eq:eis-d-warburg}、式~\eqref{eq:eis-d-cpe} 计算。}
\figsrc{结构取自 S2 §11.5.1 对四类电极 Nyquist 图与原书表 11-1 的论述
（书内 p.859—861 / PDF p.347—349）；具体数值为模型计算，非原书数据。}}
\label{fig:eis-d-model}
\end{figure}

\paragraph{读图实例：在一张模型 Nyquist 图上依次读出三类信息} 下面用图~\ref{fig:eis-d-model}
的数据走一遍完整读图流程。把\textbf{观察}与\textbf{推断}严格分开，因为这里的每一步推断都依赖
一个前提。

\begin{obs}
在图~\ref{fig:eis-d-model}(b) 的实线（含 Warburg）上（六个空心圆点即下列读数，
按频率从高到低排列）：
（1）从最高频端（80 kHz）开始，曲线左端落在 $Z'\approx 10.16$ \unit{\ohm}、
$-Z''\approx 3.97$ \unit{\ohm} 处，并未落在实轴上；
（2）频率降到 40 kHz 时 $Z'\approx 10.63$ \unit{\ohm}、$-Z''\approx 7.91$ \unit{\ohm}：
截距几乎没变，$-Z''$ 却翻了一倍；
（3）频率降到约 3.2 kHz 时，$-Z''$ 达到极大值约 50.1 \unit{\ohm}，此时
$Z'\approx 59.6$ \unit{\ohm}；
（4）继续降频，$-Z''$ 先减小到约 3.9 \unit{\ohm}（$f\approx 32$ Hz，此时
$Z'\approx 112.8$ \unit{\ohm}），再转而增大；
（5）在 $f\lesssim 10$ Hz 之后，$Z'$ 与 $-Z''$ 同步增大，在
0.1～0.02 Hz 段内 $Z'$ 从 160 \unit{\ohm} 增至 223 \unit{\ohm}、
$-Z''$ 从 50 \unit{\ohm} 增至 113 \unit{\ohm}，两点连线的斜率约为 $1.00$；
（6）把 (b) 中的点线（无扩散）与实线对比：两条线在 $f\gtrsim 100$ Hz 时肉眼不可分，
分岔从约 30 Hz 开始，到 0.02 Hz 时点线停在 $Z'=110$ \unit{\ohm}、
$-Z''\approx0$，而实线已走出 223 \unit{\ohm} 之外。
在 (a) 中对比理想弧（实线）与 CPE 弧（虚线）：两条弧的高频端都趋向同一点
（$Z'=R_{\Omega}$），低频端都趋向 $Z'=R_{\Omega}+R_{\mathrm{ct}}$（两端位置与 $n$ 无关），
但中段完全分开——理想弧顶点 $(60,\,50)$，CPE 弧顶点 $(60,\,39.4)$，后者只有前者的
$0.79$ 倍高，而顶点频率同为 $3.2$ kHz。
\end{obs}

\begin{infer}
（1）高频端不落在实轴上、且 $-Z''$ 仍有 4.0 \unit{\ohm}，说明\textbf{最高频处高频平台
尚未形成}，因此 (b) 左端那个 10.16 \unit{\ohm} \textbf{不能直接当作 $R_{\Omega}$}：
照读会把溶液电阻高估约 $1.6\%$；若改用 40 kHz 的读数 10.63 \unit{\ohm}，高估量升到
$6.3\%$。（该值是否可读，判据是 Bode 图上 $|Z|$ 与 $\varphi$ 是否出现平台。）
（2）极化电阻以弧的形式出现，说明该频段由电荷转移与双电层构成的\textbf{并联}支路支配；
在扩散尚未参与的高频段，这段弧的直径给出 $R_{\mathrm{ct}}$：理想弧的直径
$110-10=100$ \unit{\ohm}，与设定值一致。顶点条件
$\omega_{\max}=1/(R_{\mathrm{ct}}C_{\mathrm{dl}})$ 反过来给出
$C_{\mathrm{dl}}=1/(2\pi\times3183\times100)\approx0.5$ µF，也与设定值一致——
\textbf{但这两个数能读对，是因为本例恰好是完整半圆；换成 (a) 中的 CPE 弧，
用「直径」读 $R_{\mathrm{ct}}$ 仍得 100 \unit{\ohm}（两端截距与 $n$ 无关），
而用「弧顶高度 $=R_{\mathrm{ct}}/2$」读就只有 $2\times39.4=79$ \unit{\ohm}，
再代入 $C=1/(\omega_{\max}R_{\mathrm{ct}})$ 得 $0.63$ µF——与真值分别差 $21\%$ 与
$27\%$。压扁的弧不满足式~\eqref{eq:eis-d-rc} 的圆方程：顶点高度是
$(R_{\mathrm{ct}}/2)\tan(n\pi/4)$，不是 $R_{\mathrm{ct}}/2$。}
（3）$-Z''$ 先减到极小再增大，是本图上最容易被忽略的特征：极小值点的出现说明\textbf{有两个
不同时间尺度的过程串接在同一支路上}——高频弧（电荷转移）在约 30 Hz 处已衰减完毕，
而扩散项在更低频处才开始主导。因此「$-Z''$ 极小点」的位置是判断两个过程是否已分离的实用标志。
（4）低频段 $Z'$ 与 $-Z''$ 同步增大且连线斜率约 $1.00$，与式~\eqref{eq:eis-d-warburg}
的 $Z'_{\mathrm{W}}=-Z''_{\mathrm{W}}$ 一致，故该段可归为\textbf{半无限线性扩散（Warburg 阻抗）
控制}。斜率也因此是判断「扩散段是否已经独立」的实用判据：若把拟合区间一直取到弧上
（例如 3.2 kHz 与 0.02 Hz 两点连线，斜率只有约 $0.38$），正在衰减的电荷转移弧会把斜率
大幅拉低——\textbf{用斜率反算 $\sigma$ 必须先扣掉弧的贡献}。
（5）(b) 中点线与实线的分岔起点约 30 Hz，正是「扩散开始参与」的频率；而两条线在
$f\gtrsim 100$ Hz 不可分，说明\textbf{只看高频段的数据无法判断体系有没有扩散}。
这给出频率范围要求的一个定量形式：低频端必须至少比这个分岔频率低两个数量级，扩散尾才成形。
（6）(a) 中两条弧的高频端与低频端重合、中段分开，说明\textbf{弧的「两端」不携带
$n$ 的信息}——端点是 $R_{\Omega}$ 与 $R_{\Omega}+R_{\mathrm{ct}}$，与 $n$ 无关。因此
「用低频截距减高频截距得 $R_{\mathrm{ct}}$」这条做法即使在 CPE 情形下也成立
（本例中两条弧都给出 100 \unit{\ohm}）；\textbf{出错的是把中段几何（顶点高度、半高处宽度）
当作 $R_{\mathrm{ct}}/2$ 与 $1/(R_{\mathrm{ct}}C_{\mathrm{dl}})$ 去用}。
\end{infer}

\begin{pitfall}
\textbf{（1）拟合优度高不等于电路唯一。} 这是等效电路拟合的第一风险，也是本节所有其它风险的
放大器。\textbf{不同的电路可以对同一组数据给出几乎相同的拟合曲线}——只看图~\ref{fig:eis-d-model}(a)
的中段，同一段压扁的弧（跨度 $110-10=100$ \unit{\ohm}、顶点高 $39.4$ \unit{\ohm}）
至少可以用三种完全不同的电路拟合：①CPE（$n=0.85$）加 $R_{\mathrm{ct}}=100$ \unit{\ohm}；
②两个时间常数相差 $4$ 倍的 $(R\parallel C)$ 串联（各 $50$ \unit{\ohm}，即把这一段
100 \unit{\ohm} 拆成 $50+50$——峰值仍在 $3.2$ kHz、峰高 $40.0$ \unit{\ohm} 与
$39.4$ \unit{\ohm} 几乎相同，全频段与①的差别不超过约 $3$ \unit{\ohm}）；
③单个理想 $(R\parallel C)$ 串一个扩散项。三者的残差都可以做到很小，而它们给出的
「$R_{\mathrm{ct}}$」可以是 100 \unit{\ohm}（①）、两个各 50 \unit{\ohm}（②）、
或按弧顶高度误读出的 79 \unit{\ohm}——相差一倍以上。
\textbf{残差小只说明该电路在数值上能复现数据，不代表它对应真实的物理过程}。实践判据：换一组
完全不同的初值（或换一个拓扑相近的电路）重新拟合，若参数漂移到另一个解而残差几乎不变，就说明
该解不可辨识；此时只能报告「表观值」并把候选机理并列写出。

\textbf{（2）参数相关性使单个参数失去意义。} 这是风险（1）的数学根源。CPE 的 $Q$ 与 $n$ 高度
相关（$n$ 略变，$Q$ 可补偿性地变几个数量级）；$R_{\Omega}$ 与 $R_{\mathrm{ct}}$ 在同一数量级时
互相补偿；两个时间常数接近时，$R_{1}$、$C_{1}$、$R_{2}$、$C_{2}$ 四者其实只是两个可辨识的
组合量。因此\textbf{必须报告参数的置信区间与相关系数矩阵}——只报一组「最优值」而不报不确定度，
等于没有报告结果。报告 $Q$ 与 $n$ 时必须同时给出换算成电容所用的式子
（式~\eqref{eq:eis-d-cpeceff}），否则读者无法复算。

\textbf{（3）把 Warburg 当半圆。} 半无限扩散的 $45$°直线在只测到一段时很像大半径圆弧的一部分；
图~\ref{fig:eis-d-model}(b) 的实线在 10 Hz 以下那一段若单独放大，斜率 $1.00$ 的
「直线」完全可以被一个半径很大的圆拟合。直接按半圆读直径会把 $R_{\mathrm{ct}}$ 高估几倍。
\textbf{判别办法是换到 Bode 图}：扩散段的 $|Z|$ 应呈 $\omega^{-1/2}$、相角应稳定在约 $+45$°，
而圆弧段的相角会扫过一整段。

\textbf{（4）频率范围不够就下结论。} 图~\ref{fig:eis-d-model}(b) 已经给出量化例子：两条曲线
（有扩散／无扩散）在 $f\gtrsim 100$ Hz 不可分，在 40 kHz 处连高频平台都还没
形成。起始频率不够高，截距含 $R_{\mathrm{ct}}$，$R_{\Omega}$ 偏大；终止频率不够低，看不到扩散
尾或第二个弧，于是把「没看到」当成「不存在」，进而误判速控步。\textbf{频率范围不是实验参数，
它是结论的适用范围}，必须与结论一起报告。

\textbf{（5）把 CPE 当理想电容算电容值。} 这是最常见的一个算术错误，且\textbf{不报错、不告警}。
$Q$ 的量纲是 $\mathrm{S.s^{n}}$，$n=0.85$ 时 $Q=2.21\times10^{-6}$ 与同体系的
$C_{\mathrm{dl}}=0.5$ µF 相差 $4.4$ 倍，而且量纲不同——直接拿 $Q$ 当作电容会高估四倍以上；
式~\eqref{eq:eis-d-cpeceff} 的换算是必须做的。另一个同源错误是把
$C_{\mathrm{dl}}=1/(\omega_{\max}R_{\mathrm{ct}})$ 直接用在压扁的弧上——如图~\ref{fig:eis-d-model}(a)
所示，把弧顶高度当 $R_{\mathrm{ct}}/2$ 会读成 $R_{\mathrm{ct}}\approx79$ \unit{\ohm}，
代入后得 $0.63$ µF（真值 0.5 µF，高 $27\%$）；压得越扁，误差越大。
\textbf{只要 $n<0.95$，就不要用几何法读电容。}

\textbf{（6）未扣除溶液电阻。} $R_{\mathrm{ct}}$ 必须是在扣除 $R_{\Omega}$ 之后读出的弧直径；
用未校正的过电位去算动力学参数，误差随电流密度线性放大。更隐蔽的是：$R_{\Omega}$ 随电解质
浓度、电极间距与温度变化，\textbf{每次实验都应重新读，不能沿用上一批的值}。

\textbf{（7）把接触阻抗当电荷转移阻抗。} 电极与集流体、催化剂层与膜之间的接触电阻，以及参比
电极的鲁金毛细管位置不当引起的未补偿电阻，都会并入高频段，被拟合成一个大 $R_{\mathrm{ct}}$ 或
一个低 $n$ 的 CPE。这与（1）的危险叠加：接触电阻造成的低 $n$ 会被解释成「表面不均匀」这一
化学结论。\textbf{判别办法：改变电极压紧力或参比位置，若该「$R_{\mathrm{ct}}$」或「$n$」跟着
变，它就不是化学信息}。

\textbf{（8）稳态未建立就测。} 阻抗定义要求体系在测量期间是线性时不变系统。刚跳到目标电位、
表面覆盖度还在漂移、或催化剂正在失活时测得的谱，会在低频段出现「漂移型」假弧——而低频段恰好
是耗时最长、也最容易被漂移污染的一段。判据见「测量条件」第（5）条：先在单一频率上监测
$\lvert Z\rvert$ 的时间稳定性，再做全谱扫频；同一电位下正扫与反扫各测一次，两条谱重合才说明
稳态已建立。

\textbf{（9）扰动幅度过大。} 幅度过大会使响应非线性，Kramers--Kronig 检验不通过，并系统性
压低 $R_{\mathrm{ct}}$——因为大信号下电流–电势关系不再是线性的，等效电路这一整套线性工具
（包括式~\eqref{eq:eis-d-randles}）全部失效。用幅度减半检验把关。

\textbf{（10）把 $f_{\max}$ 当作 $\omega_{\max}$ 代入。} 式~\eqref{eq:eis-d-cdl} 中的
$\omega$ 是角频率。用 Hz 为单位直接代入会得到错 $2\pi$ 倍（约 $6.28$ 倍）的电容值。这是一个
纯粹的记账错误，但因为结果「看起来合理」，很少被当作可疑值复查。
\end{pitfall}

\begin{crosscheck}
阻抗谱给出的归属结论不应单独成立，至少与下列方法之一互查——每条互查针对的是\textbf{同一组
数据的一个具体参数}，而不是笼统地「再测一次」：

\textbf{（1）$R_{\mathrm{ct}}$ 对稳态极化／Tafel 斜率。} 在小过电位下
$R_{\mathrm{ct}}=RT/(nFi_0)$，因此由弧直径反推的 $i_0$ 应与 Tafel 外推值同数量级；差一个数量级
以上就说明弧里混了别的贡献（接触电阻、扩散、或第二个时间常数）。
\tocheck{未对照原页面图像：$R_{\mathrm{ct}}$ 与 $i_0$ 的换算关系式及电子数 $n$ 的约定按学科
标准形式给出，源文本 §11.5.1 未给出该式，须与页面核对。}

\textbf{（2）扩散项对旋转圆盘变转速。} 扩散控制的阻抗应在提高转速后缩短低频尾（边界层变薄），
而电荷转移控制的弧几乎不变。这条互查把 Warburg 与半圆分开，且不依赖任何模型假设——是本节最
有力的一条独立判据。RDE 的 Levich 斜率与阻抗 $\sigma$ 还可交叉校核同一个 $D$
（\S 11.4.2）。

\textbf{（3）$C_{\mathrm{dl}}$ 对循环伏安的电容电流。} 在无法拉第电流的电位区间做不同扫速的
循环伏安，由 $\mathrm{d}i/\mathrm{d}v$ 得到电容，应与阻抗高频弧给出的电容一致；两者一致才说明
高频弧确由双电层充电支配。

\textbf{（4）$R_{\Omega}$ 对电流中断法或正反馈补偿。} $R_{\Omega}$ 是阻抗谱最不可靠的一个数
（见「误判」第 6、7 条），必须用独立方法标定。三个值一致时，高频截距的归属才算落实。

\textbf{（5）时间常数个数对机理中的独立状态变量个数。} 源文本强调，从阻抗谱中时间常数的个数
及其数值大小可以获得各个步骤的动力学信息与电极表面状态变化的信息，还可观察电极过程中有无传质
过程的影响\src{书内 p.856 / PDF p.344}。把这条当作系统性的自检：\textbf{Bode 相角图上数出的
峰个数，应与所假定机理中独立状态变量的个数一致}；多了说明电路里塞进了不存在的元件，少了说明
两个过程的时间常数重叠、当前电路无法分辨它们。

\textbf{（6）与其它分片的一致性。} 本节是同一章内给出阻抗基本定义与 $Z$ 与元件关系的地方，
因此它与下一节（§11.5，电化学阻抗谱的应用与读图）的接口必须一致：\textbf{（i）}$R_{\Omega}$
的符号与含义（分片 c 沿用源文本符号 $R_{\mathrm{s}}$，指同一量）；\textbf{（ii）}Nyquist 图
纵轴取 $-Z''$ 的约定；\textbf{（iii）}相角取测量约定 $\varphi\equiv-\arg Z$（容性为正，纯电容
$+90$°、感性为负——分片 c 的读图实例即按此约定写「相角接近 $90$°」与「第四象限的负相角」）；
\textbf{（iv）}$\omega_{\max}$ 而非 $f_{\max}$ 用于
式~\eqref{eq:eis-d-cdl}。四处中任何一处不一致，读图结论就无法互相引用。

装配顺序已在 \texttt{chapters/ch11.tex} 中按文件头注释的处置办法（A）修正：本分片现排在
§11.4（分片 b）之后、§11.5 与 §11.6（分片 c）之前，故上文的「读图规则」在下文中确为
\textbf{上文}，§11.5.1 与 §11.5 的先后关系与实际排版一致。两分片之间原有一批重复段落
（阻抗定义与两种画法、等效元件表、三条读数规则、误判清单、测量条件），合并时按
\textbf{保留有源文本页码依据的一侧}处理：这些内容一律以本分片为准，分片 c 只保留
源文本 §11.5.2 特有的读图实例（甲醇氧化）与互补回查，不再重复。
\end{crosscheck}
